
We developed a practical, closed-loop elevator control pipeline that connects \emph{prediction}, \emph{traffic-mode recognition},
and \emph{proactive parking} using only standard event logs. The key contributions are:

\begin{itemize}
    \item A 5-minute, one-step-ahead demand forecasting model with online-safe features (lags, rolling statistics, and cyclical time encodings).
    \item An unsupervised traffic-mode discovery method that produces interpretable operational regimes and enables real-time classification.
    \item A mode-aware dynamic parking policy formulated as a weighted 1D $k$-median problem, yielding efficient and explainable parking targets.
    \item An interpretable benchmark and fallback (Route B) that combines a rule-based regime classifier with a state-aware zone parking allocator under rate limits, enabling auditability and safe rollback in deployment.
\end{itemize}

From an operational standpoint, improvements in \emph{tail risk} are particularly valuable: reducing P95/P99 waiting times and the long-wait
rate directly addresses passenger dissatisfaction and complaint risk. Accordingly, we emphasize a dual-route design: Route~A targets best
performance through mode-aware optimization, while Route~B provides a transparent rollback controller with rate-limited, minimal-move zone
repositioning. This combination improves both service quality and deployment readiness.

Our validation shows that the pipeline captures recurring daily regime structure and provides a deployable control rule that remains stable
under reasonable parameter perturbations. We note that the reported long-wait rate is defined relative to a configurable threshold and simulator parameterization; consequently it may be near zero in the nominal setting, while stress scenarios recover nontrivial tail risk—hence our emphasis on P95/P99 and stress-testing for decision robustness.
Future work includes integrating destination information from car calls, extending the simulator with
door times and capacity constraints calibrated to manufacturer specifications, and expanding stress-test scenarios (e.g., sharper demand spikes and mixed-regime transitions) to quantify benefits and failure modes under heavier demand, and refining online calibration of low-demand thresholds under zero-inflated counts.


